%%%%%%%%%%%%%%%%%%%%%%%%%%%%%%%%%%%%%%%%%%%%%%%%%%%%%%%%%%%%%%%%%%%%%%%%%%%%%%%
% File: ornl-template-example.tex
% Authors: Seth R Johnson, *Sam Crawford, and John Batson
% *Corresponding author, crawfordst@ornl.gov
% Date: Tuesday, February 10, 2026 @ 16:06
%
% Example of ORNL report template
%%%%%%%%%%%%%%%%%%%%%%%%%%%%%%%%%%%%%%%%%%%%%%%%%%%%%%%%%%%%%%%%%%%%%%%%%%%%%%%
%% Select the appropriate option for a TM [tm], SPR [spr], or LTR [ltr] report.
\documentclass[tm]{ornltm} % Use for TM report.
%\documentclass[spr]{ornltm} % Use for SPR report.
%\documentclass[ltr]{ornltm} % Use for LTR report.

%% Optional packages that are simply part of this example document
\usepackage{threeparttable} % Table with captions and tablenotes
\usepackage{booktabs} % Toprule, midrule, bottomrule
\usepackage{multirow} % Extend table entries over multiple rows
\usepackage{pdflscape} % Allows landscape pages for wide figures and tables

%% Bibliography styling. WARNING: the biblatex-chicago package removes the
%% footnote rule.
\usepackage[authordate,autocite=inline,backend=biber,natbib]{biblatex-chicago} % Author-date style
%\usepackage{biblatex} % Numerical style

\bibliography{example.bib} % Replace with your bib filename.

%% Use inline small italic blue for comments rather than placing them in the
%% margin.
\definecolor{note}{RGB}{48,96,192}
\renewcommand{\marginpar}[1]{{\small\em\color{note}[#1]}}

%%%%%%%%%%%%%%%%%%%%%%%%%%%%%%%%%%%%%%%%%%%%%%%%%%%%%%%%%%%%%%%%%%%%%%%%%%%%%%%
% AUTHORSHIP & PUBLISHING
%%%%%%%%%%%%%%%%%%%%%%%%%%%%%%%%%%%%%%%%%%%%%%%%%%%%%%%%%%%%%%%%%%%%%%%%%%%%%%%

%%% AUTHOR'S LIST
%% Note: multiple authors are separated with the \and command. 
%% Per ORNL guidance, please keep the number of authors at six or fewer.

\author{First Author \and Second Author \and Third Author \and Fourth Author \and Fifth Author \and Sixth Author}

%%% AUTHOR's LIST WITH AFFILIATIONS (SPR ONLY)
%% Define authors with affiliations. Does not apply to TM or LTR. Comment out if not using.
%% Note: multiple authors are separated with the \and command. 
%% Per ORNL guidance, please keep the number of authors at six or fewer.

%\author{First Author\affilnum{1} \and Second Author\affilnum{2} \and Third Author\affilnum{2} \and Fourth Author\affilnum{1} \and Fifth Author\affilnum{1} \and Sixth Author\affilnum{1}} 

%%% AFFILIATIONS (SPR ONLY)
%% Note: Does not apply to TM or LTR.

%\affiliation{\affilnum{1}Oak Ridge National Laboratory \\ \affilnum{2}University of Tennessee, Knoxville \\ }

%%% DOCUMENT TITLE
\title{Oak Ridge National Laboratory\\ Official Report Template:\\ Title Must Match RESolution Title}

%%% PUBLICATION DATE
%% Specify date or use current month and year with \today command.

%\date{April 2024}
\date{\today}

%%% REPORT NUMBER
%% Note: This number must match number provided by RESolution.

\reportnum{ORNL/XX-XXXX/XXX} % Replace first XX with TM, SPR, or LTR.

%%% DIVISION NAME
%% Define division/directorate name. Replace ``Division Name'' with your
%% division/directorate.

\division{Division Name}

%%% SPONSOR OPTIONS FOR SPR
%% Leave commented out or remove for TM and LTR.

%%% SPONSOR NUMBER
%% SPR only. Comment out if not using.

%\sponsornum{Sponsor No. (if needed)}

%%% SPONSOR LOGO(s)
%% SPR only. Add sponsor logos by listing file names below. The cover can accommodate two sponsor logos.

%\SponsorLogoOne{luggage/logo-placeholder.jpg} % Sponsor logo one. Comment out if not using.
%\SponsorLogoTwo{luggage/logo-placeholder.jpg} % Sponsor logo two. Comment out if not using.

%%%%%%%%%%%%%%%%%%%%%%%%%%%%%%%%%%%%%%%%%%%%%%%%%%%%%%%%%%%%%%%%%%%%%%%%%%%%%%%
% DOCUMENT RESTRICTIONS
%%%%%%%%%%%%%%%%%%%%%%%%%%%%%%%%%%%%%%%%%%%%%%%%%%%%%%%%%%%%%%%%%%%%%%%%%%%%%%%

%%% DRAFT
%% Mark the report as a draft with the \reportdraft command. Comment out for public release 
%% and unlimited distribution.

\reportdraft

%%% CUI
%% Adds CUI notice to the header and footer. Use % to toggle on or off. 
%% Update category, Authority Name, and Phone Number
%% For CRADA protected, use option on line 147.

%\cui{CUI//CATEGORY//DISSEM}{Authority Name, ORG}{(XXX)~XXX-XXXX} % Use for CUI basic/unspecified.
%\cui{CUI//SP-CATEGORY/SUBCATEGORY//DISSEM}{Authority Name, ORG}{(XXX)~XXX-XXXX} % Use for CUI specified.

%%% EXPORT CONTROL - Requires CUI (above)
%% Enable if your document is export controlled. Add contact info for export control authority.
%% Update Reviewer name, email, phone, date, and regulation name.

%\export{Firstname Lastname}{EmailAddress@ornl.gov}{(XXX)~XXX-XXXX}{\today}{Regulation Name}

%%%%%%%%%%%%%%%%%%%%%%%%%%%%%%%%%%%%%%%%%%%%%%%%%%%%%%%%%%%%%%%%%%%%%%%%%%%%%%%
% CRADA INSTRUCTIONS
%%%%%%%%%%%%%%%%%%%%%%%%%%%%%%%%%%%%%%%%%%%%%%%%%%%%%%%%%%%%%%%%%%%%%%%%%%%%%%%

%%% CRADA - General Info
%%
%% This command adds the CRADA notice to front cover and ORNL-851 form as the third page.
%%
%% Draft or public release is determined by toggling \reportdraft (line 99, above).
%%
%% If this report contains CRADA protected information, then toggle 
%% the \cradaprotected command (line 159, below) by uncommenting it.
%% Note that this replaces the DRAFT notice. Only use for final document.
%%

%% CRADA form ORNL-851
%%
%% Once your CRADA document is final, please fill out and sign the ORNL-851.pdf
%% form that resides in this repository. You can also download ORNL-851.pdf from
%% https://ornl.sharepoint.com/sites/forms/ornl/Forms/Forms%20View.aspx.
%%

%% Instructions for Overleaf/ShareLaTeX:
%% Download ORNL-851.pdf and open in Adobe Acrobat, fill it out, and
%% use the "Export to..." option to save the file as ORNL-851-complete.jpg.
%% Upload ORNL-851-complete.jpg to Overleaf/ShareLaTeX and replace the
%% old ORNL-851-complete.jpg file. DO NOT CHANGE THE FILE NAME. 
%% Your completed form should appear on the third page.

%% Instructions for local/repository:
%% Open ORNL-851.pdf in Adobe Acrobat, fill it out, and
%% use the "Export to..." option to save the file as ORNL-851-complete.jpg.
%% Replace old file. DO NOT CHANGE THE FILE NAME.
%% Your completed form should appear on the third page.

%\cui{CUI//SP-PROPIN/SUBCATEGORY//DISSEM}{Authority Name, ORG}{(XXX)~XXX-XXXX} % Use for CRADA protected.
%\crada{CRADA no. NFE-XX-XXXXX} % Include CRADA report number.
%\cradaprotected{ten (10) years} % If protected, include nondisclosure period in years.


%%%%%%%%%%%%%%%%%%%%%%%%%%%%%%%%%%%%%%%%%%%%%%%%%%%%%%%%%%%%%%%%%%%%%%%%%%%%%%%
% ABBREVIATIONS
%%%%%%%%%%%%%%%%%%%%%%%%%%%%%%%%%%%%%%%%%%%%%%%%%%%%%%%%%%%%%%%%%%%%%%%%%%%%%%%
%% Uses the Acro package called in the CLS (line 61) and styled in the STY (lines
%% 135--144). Note that acronyms that aren't used in the document won't be printed
%% in the list. Declare your abbreviations in abbreviations.tex.

%%%%%%%%%%%%%%%%%%%%%%%%%%%%%%%%%%%%%%%%%%%%%%%%%%%%%%%%%%%%%%%%%%%%%%%%%%%%%%%
% ABBREVIATIONS
%%%%%%%%%%%%%%%%%%%%%%%%%%%%%%%%%%%%%%%%%%%%%%%%%%%%%%%%%%%%%%%%%%%%%%%%%%%%%%%
%% Uses the Acro package called in CLS and styled in the STY. Note that acronyms
%% that aren't used in the document won't be printed in the list.
%% Declare your abbreviations below.

\DeclareAcronym{doe}{
  short = DOE ,
  long = US Department of Energy
}

\DeclareAcronym{nrc}{
  short = NRC ,
  long = Nuclear Regulatory Commission
}

\DeclareAcronym{ornl}{
  short = ORNL ,
  long = Oak Ridge National Laboratory
}

\DeclareAcronym{tnr}{
  short = TNR ,
  long = Times New Roman
}

\DeclareAcronym{tm}{
  short = TM ,
  long = Technical Memo
}

%%%%%%%%%%%%%%%%%%%%%%%%%%%%%%%%%%%%%%%%%%%%%%%%%%%%%%%%%%%%%%%%%%%%%%%%%%%%%%%
% FRONT MATTER
%%%%%%%%%%%%%%%%%%%%%%%%%%%%%%%%%%%%%%%%%%%%%%%%%%%%%%%%%%%%%%%%%%%%%%%%%%%%%%%
\begin{document}
\frontmatter

% Prints TABLE OF CONTENTS
\tableofcontents % Comment out for LTR

% Prints LIST OF FIGURES
\listoffigures % Comment out for LTR

% Prints LIST OF TABLES
\listoftables % Comment out for LTR

%% Prints LIST OF ABBREVIATIONS
\printacronyms % Comment out for LTR. Requires acro package.

%%%%%%%%%%%%%%%%%%%%%%%%%%%%%%
%%% ADDITIONAL FRONT MATTER (OPTIONAL)
\begin{foreword} % Remove foreword environment if not using (lines 217--291).

Material after the Table of Contents, List of Figures, List of Tables, and List of
Abbreviations may include any or all of the following sections and should begin
on an odd-numbered page in the order listed below:

\begin{itemize}
  \item Foreword \marginpar{Note: spelling is not ``Forward.''}
  \item Preface
  \item Acknowledgments
  \item Executive summary, or Summary
  \item Abstract
\end{itemize}

Each section of the front material begins on an odd-numbered page (the
\verb|\section| command inserts a \verb|\cleardoublepage| before each one). The
abstract, if very brief, may begin on page 1 inserted above the introduction
section.

%%%%%%%%%%%%%%%%%%%%%%%%%%%%%%
\subsection*{TM, SPR, or LTR?}

\textbf{Technical Memos (TMs)} are the primary report type produced at the lab,
and TM is the default option for this document's class. The TM template has been
approved by DOE for ORNL's use. The picture on the front cover is the only
design or branding element that may be changed. Use the RESolution Publications
system to generate an ORNL/TM report number by entering the report in the system
and selecting the TM category.

\textbf{Letter reports (LTRs)} are for externally published letter reports, not
letters or sponsor reports. Use the RESolution Publications system to generate
an ORNL/LTR report number by entering the report in the system and selecting the
LTR category. Notably, LTR reports should be limited to 10 pages of narrative
text. Longer documents should use the TM option.

\textbf{Sponsor reports (SPRs)} are for reports funded by outside sponsors that
have not provided their own guidelines and/or templates. If the sponsor (e.g.,
Nuclear Regulatory Commission) has provided a template, then use it. Otherwise,
use the ORNL SPR option in this example. Notably, even if you are using a
sponsor's template, you must still use the RESolution Publications system to
generate an ORNL/SPR report number by entering the report in the system and
selecting the SPR category.

%%%%%%%%%%%%%%%%%%%%%%%%%%%%%%
\subsection*{General Information}

Document margins are 1 in.~all around. Generally, the \ac{tm} styling reflects
the defaults of Microsoft Word because that is the publishing application of
choice for the \ac{doe}. \Acp{tm} are more fun to write in \LaTeX.

%%%%%%%%%%%%%%%%%%%%%%%%%%%%%%
\subsection*{Footnotes}

Text footnotes are 9 pt. \ac{tnr} with no space between notes.\footnote{Footnotes
are 9 pt. \ac{tnr} with no space between notes.} Footnote indicators in text can be
superscript Arabic numerals as shown above or as superscript symbols
(using markup changes) as shown below.

%% Use symbols for footnotes
\renewcommand*{\thefootnote}{\fnsymbol{footnote}}

If symbols, then they will appear in the following order: \footnote{First},
\footnote{Second}, \footnote{Third}, \footnote{Fourth}, \footnote{Fifth}, and
\footnote{etc}. Symbols \emph{are not} used as table footnote indicators (see
Section~\ref{sec:tables}).

%% Use Arabic numbers for footnotes
\renewcommand*{\thefootnote}{\arabic{footnote}}
% Restore footnote count
\setcounter{footnote}{0}

Footnotes may be used in addition to the author-date citation method, in which
case, the footnotes provide additional information to the text---not
bibliographic citations.
\end{foreword}

%%%%%%%%%%%%%%%%%%%%%%%%%%%%%%%%%%%%%%%%%%%%%%%%%%%%%%%%%%%%%%%%%%%%%%%%%%%%%%%
% ABSTRACT
%%%%%%%%%%%%%%%%%%%%%%%%%%%%%%%%%%%%%%%%%%%%%%%%%%%%%%%%%%%%%%%%%%%%%%%%%%%%%%%
\mainmatter
\begin{abstract}

The abstract, if short, is inserted here (after \verb|\mainmatter|) before the
introduction section. If the abstract is long, then it should be inserted
\emph{before} \verb|\mainmatter|, and the introduction below will be the first
entry on the page. This \acl{tm} class is used for reports at \acl{ornl}.

\end{abstract}

%%%%%%%%%%%%%%%%%%%%%%%%%%%%%%%%%%%%%%%%%%%%%%%%%%%%%%%%%%%%%%%%%%%%%%%%%%%%%%%
% DOCUMENT CONTENT
%%%%%%%%%%%%%%%%%%%%%%%%%%%%%%%%%%%%%%%%%%%%%%%%%%%%%%%%%%%%%%%%%%%%%%%%%%%%%%%
\acresetall % Reset abbreviations for narrative.
\section{FIRST-ORDER HEADING}
\label{sec:firstsec}
\marginpar{Sections and subsections must be manually capitalized to render
properly in the table of contents.}

Begin text here. You can reference abbreviations such as \ac{ornl}, and, because
we are using the \verb|Acro| package, they will be defined on first instance and
automatically shortened on subsequent appearances---e.g., \ac{ornl}---unless you
reset the counter by using \verb|\acresetall|. They will also appear in the List
of Abbreviations. Note that any abbreviations used prior to the narrative are
reset! They are reset again for the appendix because these elements stand alone.

The body of the document is paginated consecutively (i.e., 1, 2, 3) or by
section (i.e., 1-1, 2-1, 3-1) using Arabic numbers centered at the bottom of
each page. Blank backup pages are not numbered.

When typing your document in LaTeX, it might be useful to break the line at the
right-hand side by hitting the return key so that any version control system you
might be using can recognize changes with better granularity. These version
control systems (e.g., SVN, Git, Mercurial) track changes by line number, so the
more frequent line breaks that a document has, the better control over changes
and merges the author(s) will have. These breaks are treated as spaces in the
final output. We recommend breaking lines at 80 characters. Some text editors
have a hotkey and/or extensions for this procedure (e.g., Rewrap for VS Code).

%%%%%%%%%%%%%%%%%%%%%%%%%%%%%%%%%%%%%%%%%%%%%%%%%%%%%%%%%%%%%%%%%%%%%%%%%%%%%%%
\subsection{SECOND-ORDER HEADING}

Begin text here.

Another paragraph is here.

\subsubsection{Third-Order Heading}

Begin text.

\paragraph{Fourth-Order Heading}

Begin text here.

\subparagraph{Fifth-Order Heading}

Ask yourself, ``Do I \emph{really} need five or more hierarchical levels in my
document?''


%%%%%%%%%%%%%%%%%%%%%%%%%%%%%%%%%%%%%%%%%%%%%%%%%%%%%%%%%%%%%%%%%%%%%%%%%%%%%%%
\marginpar{A \texttt{\textbackslash cleardoublepage} should be inserted here if the preceding
section is not too short.}

\section{FIGURE PLACEMENT}
\label{sec:figures}

\subsection{FIGURES}

Insert figures after the text callout but as close to the callout as possible.
They always have a caption describing them, and they are always numbered. LaTeX
automatically floats tables and figures depending on how much space is left on
the page where they are processed. If there is not enough room on
the current page, then the float is moved to the top of the next page. Figures are
numbered consecutively (i.e., Figure~1, Figure~2) except in large reports, in
which they may be numbered consecutively by section (i.e., Figure~1.1,
Figure~1.2).

All figures are pulled in from source files that should be placed in a directory
that the author can link to from the LaTeX file.  In this template, the image
files for the figures are in a subfolder named ``figures,'' which is
standard practice for LaTeX documents.

\begin{figure}[ht]
  \centering
	\includegraphics[width=5.5in, height=2.5in]{figures/fig1.png}
  \longcaption{All figure captions are 10 pt, bold, \ac{tnr}.}{Only the first sentence of
  a figure caption is bold.}
\end{figure}

Text can go in between figures. LaTeX will pretty much place the figures
wherever it sees fit based on the parameters described above---unless you
explicitly command it to do otherwise.

\begin{figure}[ht]
  \centering
	\includegraphics[width=6.2in, height=2.75in, keepaspectratio=true]{figures/fig1.png}
	\caption{One-line figure captions are bold.}
\end{figure}

%%%%%%%%%%%%%%%%%%%%%%%%%%%%%%%%%%%%%%%%%%%%%%%%%%%%%%%%%%%%%%%%%%%%%%%%%%%%%%%

\section{TABLE PLACEMENT}
\label{sec:tables}

\subsection{TABLES}

As with figures, tables are numbered consecutively or consecutively by section
and should follow their callout (see Table \ref{tab:1}) in the text as closely
as possible. LaTeX automatically floats tables and figures depending on how
much space is left on the page where they are processed. If there is
not enough room on the current page, then the float is moved to the top of the next
page. Tables are formatted using small \ac{tnr}. To accommodate large tables, the
font size can be decreased and, if necessary, landscaped (Table~\ref{tab:2}).

\begin{table}[ht]
  \centering
  \begin{threeparttable}
  \caption{Table caption is bold, centered, and initial cap with no period at
  end of title}\label{tab:1}
  \newcommand\tworow[1]{\multirow{2}{*}{#1}}
  \begin{tabular}{c c c c c c}
    \toprule
    \tworow{EM project\tnote{a}} & \tworow{Recycling method} & \tworow{Amount
    recycled (lb)} & \multicolumn{3}{c}{Cost (\$)} \\ \cline{4-6}
    &  &  & Recycling & Disposal &  Storage \\ \midrule
    Metals recycle & Smelting & 1,072,000 & 1,565,763 & 1,338,447 & 1,608,000 \\
    Cooling tower\tnote{b} & Decontamination & 459,000 & 605,880 & 573,120 &
    688,500  \\
    \em{Total} &   & 1,601,150 & 2,266,491 & 2,004,973 & 2,401,725 \\
    \bottomrule
  \end{tabular}
  \begin{tablenotes}
    \item \emph{Note:} A general note to the table as a whole is not linked to
    a superscript letter. It is formatted like this note. Table footnotes are
    9~pt, or, as with larger tables, 1~pt smaller than the table body text.
  \item[a] Footnote callouts are lowercase, italic, superscript letters in
    sequence from left to right.
  \item[b] For multipage tables, the footnote appears only on the last page of
    the table.
  \end{tablenotes}
  \end{threeparttable}
\end{table}

\begin{landscape} % Creates a page break and orients the new page(s) in landscape mode, including the header and footer.
\begin{table}
  \caption{This table is too wide for a portrait orientation and uses the pdflscape package for landscape}\label{tab:2}
  \centering
  \begin{tabular}{llllllllllll}
  \toprule
  Key  & Def.     & Definition           & Key  & Def. & Definition            & Key  & Def. & Definition           \\
  \midrule
  RND= & given    & random number        & PNU= & NO   & use delayed neutron   & SMU= & NO   & self-multiplication  \\
  TME= & no limit & execution time (min) & AMX= & NO   & all mixture xsecs     & NUB= & YES  & neutrons per fission \\
  TBA= & 10 min   & batch time (min)     & XAP= & NO   & xsec angles \& probs. & PAX= & NO   & albedo-xsec array    \\
  WTA= & 0.5      & average weight       & XS1= & NO   & 1D xsecs             & TFM= & NO   & coordinate transform \\
  WTH= & 3.0      & wt. for splitting    & XS2= & NO   & 2D xsecs             & PMF= & NO   & print angular fluxes \\
  WTL= & 1/WTH    & Russian Roulette wt. & XSL= & NO   & 2D Pl xsecs          & MFX= & NO   & compute mesh fluxes  \\
  SIG= & 0.0      & deviation limit      & PKI= & NO   & fission spectrum      & PMS= & NO   & print mesh fluxes    \\
  \bottomrule
  \end{tabular}
\end{table}
\end{landscape}

\section{FIGURE WRAPPING}
\label{sec:figwrap}

The template uses the \verb|wrapfig| package to enable text to wrap neatly
around figures. You can adjust the padding and placement by using the
\verb|trim| and \verb|width| variables.

Lorem ipsum dolor sit amet, consectetur adipiscing elit. Etiam elementum euismod
rutrum. Cum sociis natoque penatibus et magnis dis parturient montes, nascetur
ridiculus mus. Suspendisse gravida nunc est, at viverra augue rhoncus et. In
varius luctus enim, venenatis ullamcorper ipsum volutpat at. Pellentesque
habitant morbi tristique senectus et netus et malesuada fames ac turpis egestas.
Aliquam vel blandit sem. Etiam vel erat nibh. Vestibulum vulputate, neque vel
aliquet consequat, tortor magna sagittis magna, nec bibendum nulla ligula ut
nunc. Curabitur sed lorem in orci vestibulum imperdiet.

\begin{wrapfigure}{r}{7cm} % Set box alignment {r} and size {7cm}
  \centering % Center figure within box
    \includegraphics[trim=2cm 0cm 2cm 0.5cm,width=4.5cm]{example-image-a} %Set image dimension smaller than box
  \caption{Wrapped figure.} % Add caption
  \vspace{-0.5cm} % Adjust white space under caption.
\end{wrapfigure}

Curabitur eleifend vel leo eu malesuada. Donec ultricies cursus ante, sit amet
tincidunt enim ultrices id. Sed vel odio vel risus gravida aliquam vel eu neque.
Donec in ipsum ac risus varius aliquam facilisis sed neque. Nunc suscipit neque
at augue volutpat viverra. Proin sit amet nunc a lectus condimentum elementum.
Nullam porttitor felis ipsum, sit amet facilisis felis volutpat et.

Quisque bibendum ligula ut sapien placerat, vitae faucibus neque luctus.
Maecenas semper elementum tortor id laoreet. Aenean condimentum tellus ante, et
condimentum magna sodales quis. Morbi interdum lobortis massa at pulvinar. Proin
et suscipit mauris. Maecenas aliquam lorem ipsum, nec dapibus nibh aliquam
vitae. Sed et molestie sapien. Nullam commodo id mauris quis mollis.

%%%%%%%%%%%%%%%%%%%%%%%%%%%%%%%%%%%%%%%%%%%%%%%%%%%%%%%%%%%%%%%%%%%%%%%%%%%%%%%

\section{CITATIONS}
\label{sec:citations}

\subsection{USING IN-TEXT CITATIONS}
For the \ac{ornl} \ac{tm} reports, \textit{Chicago Manual of Style} references and
citations are preferred. Of these, you can choose between author-date +
references cited or footnote + references cited.

Author-date style is set by default and is preferred. Use
the \verb|\autocite| command to defer to the default citation style, which will give
you the following: \autocite{knuthwebsite}.

If for some reason you don't want the citation to be surrounded by parentheses,
you can use the \verb|\cite| command like this: \cite{einstein}. Or you can place
parentheses around the year by using the \verb|\citet| command like this:
\citet{einstein}.

You can also push abbreviated citations into the footnotes using a \verb|\footcite|
command, like this. \footcite{latexcompanion}

For all of the examples above, you do not have to change anything other than the
actual citation command. However, you (advanced users) can change the
bibliography further by changing the bibliography styles on line 23.

%%%%%%%%%%%%%%%%%%%%%%%%%%%%%%%%%%%%%%%%%%%%%%%%%%%%%%%%%%%%%%%%%%%%%%%%%%%%%%%

\section{LISTS}
\label{sec:lists}

Flanges used in ultra-high vacuum service are commonly closed with metallic
gaskets or seals. The geometry of the seal interface is critical for proper
function. Experience has shown that flanges subjected to conventional welding
techniques distort beyond a point that they can be sealed.
\vspace{12pt}

\begin{itemize}
  \item \textbf{Apply itemize to format an unordered list with bullets.}
    \emph{Add extra space between items if preferred.} Flanges used in
    ultra-high vacuum service are commonly closed with metallic gaskets or
    seals. The geometry of the seal interface is critical for proper function.
    Experience has shown that flanges subjected to conventional welding
    techniques distort beyond a point that they can be sealed.
	\begin{itemize}
    \item \textbf{Apply second itemize to format a sublist.} Flanges
      used in ultra-high vacuum service are commonly closed with metallic
      gaskets or seals. The geometry of the seal interface is critical for
      proper function. Experience has shown that flanges subjected to
      conventional welding techniques distort beyond a point that they can be
      sealed.
	\end{itemize}
  \item \textbf{Apply itemize for format ``Bullet (last item)'' to format the
    ending item.} Flanges used in ultra-high vacuum service are commonly closed
    with metallic gaskets or seals. The geometry of the seal interface is
    critical for proper function. Experience has shown that flanges subjected to
    conventional welding techniques distort beyond a point that they can be
    sealed.
\end{itemize}

\subsection{SECOND-ORDER HEADING}

Flanges used in ultra-high vacuum service are commonly closed with metallic
gaskets or seals. The geometry of the seal interface is critical for proper
function. Experience has shown that flanges subjected to conventional welding
techniques distort beyond a point that they can be sealed.

You can also create an ordered/numbered list with \verb|enumerate|.

\begin{enumerate}
  \item Flanges used in ultra-high vacuum service are commonly closed with
    metallic gaskets or seals. The geometry of the seal interface is critical
    for proper function. Experience has shown that flanges subjected to
    conventional welding techniques distort beyond a point that they can be
    sealed.
  \item Flanges used in ultra-high vacuum service are commonly closed with
    metallic gaskets or seals. The geometry of the seal interface is critical
    for proper function. Experience has shown that flanges subjected to
    conventional welding techniques distort beyond a point that they can be
    sealed.
\end{enumerate}

\subsubsection{Third-Order Heading}

Lorem ipsum dolor sit amet, consectetur adipiscing elit.

\paragraph{Fourth-order heading}

Cras elementum pellentesque lacinia. Pellentesque habitant morbi tristique
senectus et netus et malesuada fames ac turpis egestas. Curabitur ornare libero
ipsum, quis volutpat diam sagittis porta. Aenean rhoncus dignissim leo, sed
lobortis lectus tempus venenatis. Proin pulvinar lectus a tellus finibus, quis
sodales nisl ullamcorper. Aliquam vel ornare ex, tristique faucibus velit. Fusce
scelerisque, lacus eget tincidunt varius, quam diam bibendum nulla, eu molestie
urna neque ac eros. Vestibulum dolor lectus, bibendum a mi nec, rhoncus
imperdiet ante. Cum sociis natoque penatibus et magnis dis parturient montes,
nascetur ridiculus mus. Duis gravida orci nec bibendum lobortis. Morbi leo
magna, ultricies eu interdum id, volutpat ac justo. Quisque mattis, turpis eu
interdum congue, augue odio ultricies ipsum, quis venenatis magna diam a nisl.
Vivamus finibus, mi vitae rhoncus pretium, mauris turpis fermentum dolor, id
volutpat nunc mauris sed urna.

\paragraph{Fourth-order heading}

Aliquam volutpat maximus orci nec dignissim. Sed varius, tellus non auctor
convallis, lorem enim convallis dolor, at ultrices libero tortor rutrum arcu.
Cras imperdiet sit amet nisl quis iaculis. Sed et tempus libero. Aenean a nisl
metus. Etiam ut orci id libero pretium pretium. Nam mauris est, pretium lobortis
leo ut, scelerisque posuere sapien. Vivamus posuere venenatis sem at
pellentesque.

\paragraph{Fourth-order heading}

Aenean imperdiet vestibulum faucibus. Aliquam id nulla elementum ipsum
scelerisque viverra. Pellentesque fringilla ultricies libero interdum dictum.
Quisque id pellentesque nisl, posuere convallis urna. Aliquam urna neque,
molestie efficitur congue nec, ornare vitae erat. Duis a sagittis ante. Aliquam
lacinia placerat dapibus. In in finibus dolor. Morbi id aliquet ante, et mattis
quam. Nunc rutrum viverra leo, et eleifend nulla interdum vitae. Donec nec
finibus nibh. Integer interdum nibh enim, id porta libero vestibulum vel. Sed
efficitur justo sed diam suscipit varius.

%%%%%%%%%%%%%%%%%%%%%%%%%%%%%%%%%%%%%%%%%%%%%%%%%%%%%%%%%%%%%%%%%%%%%%%%%%%%%%%
% REFERENCES
%%%%%%%%%%%%%%%%%%%%%%%%%%%%%%%%%%%%%%%%%%%%%%%%%%%%%%%%%%%%%%%%%%%%%%%%%%%%%%%

% Bibliography support via biblatex
\section{REFERENCES}% Create heading for references to include in ToC.
\printbibliography[heading=none]% Print references without biblatex's heading.

%%%%%%%%%%%%%%%%%%%%%%%%%%%%%%%%%%%%%%%%%%%%%%%%%%%%%%%%%%%%%%%%%%%%%%%%%%%%%%%
% APPENDICES
%%%%%%%%%%%%%%%%%%%%%%%%%%%%%%%%%%%%%%%%%%%%%%%%%%%%%%%%%%%%%%%%%%%%%%%%%%%%%%%
\appendix
\acresetall % Reset abbreviations for Appendices.
\section{FIRST APPENDIX}
\label{sec:firstapp}

If each appendix contains similarly formatted text as the body of the document
(i.e., first-order headings), then flysheets are not necessary. If first-order
headings cannot be used (e.g., with computer data or forms), then a flysheet should
be used.

\begin{equation}\label{eq:appeq}
  2 + 2 = 4
\end{equation}

\begin{figure}[htb]
  \centering
	\includegraphics[width=5.5in, height=2.5in]{example-image-a}
  \caption{Figure caption.}
\end{figure}

%%%%%%%%%%%%%%%%%%%%%%%%%%%%%%%%%%%%%%%%%%%%%%%%%%%%%%%%%%%%%%%%%%%%%%%%%%%%%%%
\section{SECOND APPENDIX}
\label{sec:secondapp}
If each appendix contains similarly formatted text as the body of the document
(i.e., first-order headings), then flysheets are not necessary. If first-order
headings cannot be used (e.g., with computer data or forms), then a flysheet should
be used.

\begin{table}[h]
  \centering
  \caption{A regular table in the appendix}
  \label{tab:app-table}
  \begin{tabular}{ccccc}
  \toprule
  Class & Precision & Recall & F1-Score & Support \\
  \midrule
  0 & 1.00 & 1.00 & 1.00 & 476075 \\ 
  1 & 1.00 & 0.99 & 0.99 & 81294 \\ \midrule
  Accuracy & - & - & 1.00 & 557369 \\ 
  Macro Avg. & 1.00 & 0.99 & 1.00 & 557369 \\ 
  Weighted Avg. & 1.00 & 1.00 & 1.00 & 557369 \\ \bottomrule
  \end{tabular}
  \end{table}


%%%%%%%%%%%%%%%%%%%%%%%%%%%%%%%%%%%%%%%%%%%%%%%%%%%%%%%%%%%%%%%%%%%%%%%%%%%%%%%
% Back cover page
\backmatter

%%%%%%%%%%%%%%%%%%%%%%%%%%%%%%%%%%%%%%%%%%%%%%%%%%%%%%%%%%%%%%%%%%%%%%%%%%%%%%%
\end{document}
%%%%%%%%%%%%%%%%%%%%%%%%%%%%%%%%%%%%%%%%%%%%%%%%%%%%%%%%%%%%%%%%%%%%%%%%%%%%%%%
% end of ornltm/ornl-template-example.tex
%%%%%%%%%%%%%%%%%%%%%%%%%%%%%%%%%%%%%%%%%%%%%%%%%%%%%%%%%%%%%%%%%%%%%%%%%%%%%%%
